
\section{Introduction}

An SCC application can be defined according to an architectural
pattern involving four kinds of components, organized into layers: (1)
\emph{sensors} at the bottom, which obtain information about the
environment; (2) then \emph{context operators}, which process this
information; (3) then \emph{control operators}, which use this refined
information to control (4) \emph{actuators} at the top, which finally
impact the environment.

\begin{figure}
  \centering
  \includegraphics[width=\linewidth]{diaspec-process}
  \caption{Overview of the \diaspec{} development process}
  \label{fig:diaspec-process}
\end{figure}

\diaspec{} is a design language dedicated to describing SCC
applications. From such a design description, the \diaspec{} compiler
produces a dedicated Java programming framework
(Figure~\ref{fig:diaspec-process}) that is both \emph{prescriptive}
and \emph{restrictive}: it is prescriptive in the sense that it guides
the developer, and it is restrictive in the sense that it limits the
developer to what the design description allows. For each component
description in the design, the compiler generates an abstract class.
The abstract methods in this class represent code to be provided by
the developer, to allow him to program the application logic. To
implement a component, the developer implements these methods in a
subclass of this abstract class and thus never modifies generated
code.

By separating application logic (implemented by the developers) and
runtime support (generated in the programming framework), \diaspec{}
facilitates the design, implementation and evolution of SCC
applications. In the rest of this paper, we assume that robotics
applications are a subset of SCC applications and demonstrate that
\diaspec{} can be applied in the robotics domain.

%%% Local Variables:
%%% mode: latex
%%% coding: utf-8
%%% TeX-master: "dslrob"
%%% TeX-PDF-mode: t
%%% ispell-local-dictionary: "english"
%%% End:


%%% Local Variables: 
%%% mode: latex
%%% TeX-master: "dslrob"
%%% End: 
